%!TEX root =  tfg.tex
\chapter{Arranque}

\begin{quotation}[Novelist]{Ernest Hemingway (1899--1961)}
The good parts of a book may be only something a writer is lucky enough to overhear or it may be the wreck of his whole damn life -- and one is as good as the other.
\end{quotation}

\begin{abstract}
En este capítulo se formaliza el alcance del proyecto y se 
establecen las bases técnicas para su desarrollo. El objetivo es 
traducir la visión conceptual de Via Inferni en una lista estructurada de 
requisitos funcionales, aplicando una primera iteración de la metodología Feature 
Driven Development (FDD). Asimismo, se define la arquitectura de software que soportará
dichas funcionalidades, detallando los entornos de producción, las herramientas de 
diseño de datos (preproducción) y la estrategia de pruebas necesaria para asegurar 
la calidad del software.
\end{abstract}

\section{Lista de características}

Siguiendo la metodología \textit{Feature Driven Development} (FDD), 
se ha descompuesto la visión global del videojuego en conjuntos de características 
agrupados por dominio funcional. Esta lista representa el alcance del 
Producto Mínimo Viable (MVP) y prioriza las mecánicas 
centrales que definen la experiencia \textit{roguelike} en la ambientación de 
la Divina Comedia. El objetivo de esta sección es describir, de forma clara y 
comprensible, qué elementos formarán parte de la primera versión funcional del juego 
y qué papel cumple cada uno dentro de la experiencia completa.

\subsection{Dominio: Jugador y Mecánicas Core}
Este conjunto reúne las funcionalidades esenciales del jugador, es decir, aquellas 
sin las cuales no sería posible jugar una partida completa. Incluye el control del 
personaje, la interacción con los enemigos y las reglas básicas que determinan cuándo 
una partida continúa y cuándo termina.
\begin{itemize}
    \item \textbf{Gestión de Personajes:} Se implementará un sistema de dos personajes 
    jugables, con capacidad para distinguirlos y gestionar cuál está activo en cada 
    momento de la partida.
    \begin{itemize}
        \item La idea es que ambos personajes compartan la base del control, pero puedan 
        tener diferencias en atributos o comportamiento para aportar variedad jugable.
    \end{itemize}
    \item \textbf{Sistema de Combate:}
    \begin{itemize}
        \item Se desarrollará el sistema que permite atacar, recibir daño y resolver los 
        enfrentamientos contra enemigos y jefes.
        \item Este sistema incluirá la lógica básica de impacto, daño y respuesta visual 
        para que el combate resulte comprensible para el jugador.
    \end{itemize}
    \item \textbf{Ciclo de Vida (Run):}
    \begin{itemize}
        \item Se controlará la salud del jugador, el estado de muerte y la condición que 
        determina el final de la partida.
        \item Cuando el personaje muera, la partida se reiniciará para mantener la 
        estructura típica de un roguelike con progresión por intentos.
    \end{itemize}
\end{itemize}

\subsection{Dominio: Generación Procedural y Mundo}
Este bloque describe cómo se construirá el escenario de forma automática en cada partida. 
En lugar de usar un único mapa fijo, el juego generará niveles diferentes dentro de los 
nueve círculos del infierno, lo que aporta rejugabilidad y hace que cada intento sea distinto.
\begin{itemize}
    \item \textbf{Algoritmo de Generación:}
    \begin{itemize}
        \item Se generará una estructura de habitaciones conectadas entre sí siguiendo una 
        lógica procedimental, de manera que el mapa cambie en cada partida.
        \item Las salas no serán todas iguales: se escogerán distintos \textit{layouts} 
        prediseñados según el tipo de habitación, como salas de combate, tesoro, tienda o descanso.
        \item La dificultad y la complejidad del mapa aumentarán progresivamente a medida que 
        el jugador avance por círculos más profundos.
    \end{itemize}
    \item \textbf{Navegación:}
    \begin{itemize}
        \item Se controlará el estado de las puertas para evitar que el jugador avance mientras 
        queden enemigos vivos en una sala.
        \item El paso de un círculo al siguiente se activará al derrotar al jefe correspondiente.
    \end{itemize}
\end{itemize}

\subsection{Dominio: Entidades y Comportamiento}
Este apartado incluye a los enemigos y jefes que el jugador encontrará durante la partida. 
Su función es generar desafío, obligar al jugador a adaptar su forma de jugar y dar identidad 
propia a cada zona del juego.
\begin{itemize}
    \item \textbf{Arquetipos de Enemigos:} Cada uno de los nueve círculos contará con tres 
    tipos de enemigo claramente definidos, de forma que la dificultad y la variedad aumenten 
    de manera controlada a lo largo de la partida.
    \begin{itemize}
        \item \textbf{Melee:} Enemigos de corto alcance con un comportamiento estándar, 
        pensados para constituir la base del combate en cada círculo.
        \item \textbf{Flying:} Enemigos que se desplazan por el escenario evitando parte de 
        las limitaciones del terreno, lo que los hace más difíciles de controlar.
        \item \textbf{Tank:} Enemigos más resistentes y peligrosos, diseñados para resistir 
        más daño y obligar al jugador a prestarles una atención especial.
    \end{itemize}
    \item \textbf{Jefes Finales:}
    \begin{itemize}
        \item \textbf{Cancerbero:} Jefe de un nivel avanzado, con un conjunto de ataques y 
        patrones específicos que funcionarán como prueba de habilidad para el jugador.
        \item \textbf{Satán:} Jefe final del recorrido, diseñado como el mayor reto de la partida 
        y encargado de cerrar el progreso del jugador dentro del juego.
    \end{itemize}
\end{itemize}

\subsection{Dominio: Economía y Progresión}
En este bloque se agrupan los sistemas que hacen que la partida no se base solo en combatir, 
sino también en tomar decisiones sobre recursos y mejoras. Estos elementos dan sentido al 
progreso del jugador dentro de una misma run.
\begin{itemize}
    \item \textbf{Sistema de Ítems:}
    \begin{itemize}
        \item Se implementarán objetos que modifiquen las estadísticas del personaje de forma 
        permanente o mientras dure la partida.
        \item También habrá consumibles con efectos puntuales, pensados para resolver situaciones 
        concretas o dar una ventaja temporal.
    \end{itemize}
    \item \textbf{Economía:}
    \begin{itemize}
        \item Los enemigos podrán soltar monedas u otros recursos al ser derrotados, para que el 
        jugador obtenga recompensas por avanzar y combatir.
        \item Esos recursos se podrán gastar en la sala de tienda para adquirir mejoras, objetos 
        o ventajas que ayuden a progresar.
    \end{itemize}
\end{itemize}

\subsection{Dominio: Interfaz y Feedback (UI/UX)}
\begin{itemize}
    \item \textbf{HUD (Heads-Up Display):} Se mostrará en pantalla la información 
    imprescindible para jugar, como la vida, el personaje activo, el minimapa y los tiempos 
    de espera de habilidades o acciones.
    \item \textbf{Flujo de Pantallas:} Se implementará la navegación entre las pantallas 
    principales del juego, incluyendo el menú inicial, la pausa durante la partida y la 
    pantalla final de resultados.
\end{itemize}

\section{Diseño arquitectónico}

\subsection{Visión general del sistema}

El videojuego desarrollado en el presente Trabajo de Fin de Grado consiste en un roguelike con generación procedural de niveles inspirado en la estructura del Infierno de Dante. El juego se organiza en nueve círculos temáticos, cada uno con identidad visual propia, enemigos específicos, conjunto de salas diferenciado y un jefe final que permite el acceso al siguiente círculo.

Desde el punto de vista arquitectónico, el sistema se ha diseñado siguiendo una estructura modular orientada a la separación de responsabilidades y al bajo acoplamiento entre componentes. La arquitectura se basa en un conjunto de gestores especializados (manager-based architecture), donde cada subsistema tiene una responsabilidad claramente delimitada.

Adicionalmente, el sistema contempla la existencia de dos personajes jugables, Dante y Virgilio, cuya lógica de control y estado debe poder alternarse dinámicamente durante la ejecución. Aunque la mecánica de diferenciación entre ambos personajes no se encuentra completamente definida en esta fase, la arquitectura global se ha diseñado para permitir su extensión sin modificaciones estructurales significativas.

\subsection{Arquitectura global del sistema}

El sistema se organiza en los siguientes subsistemas principales:

\begin{itemize}
    \item \textbf{GameManager}: Controlador global del estado del juego.
    \item \textbf{CircleManager}: Gestión del círculo actual.
    \item \textbf{LevelGenerator}: Generación estructural del nivel.
    \item \textbf{RoomManager}: Control de activación de salas.
    \item \textbf{RoomGenerator}: Generación de contenido interno.
    \item \textbf{DifficultyManager}: Ajuste dinámico de dificultad.
    \item \textbf{CharacterManager}: Gestión de personajes jugables.
\end{itemize}

El siguiente diagrama ilustra la estructura general y las relaciones entre subsistemas:

\begin{center}
\begin{tikzpicture}[
    node distance=1.8cm,
    box/.style={rectangle, draw, rounded corners, align=center, minimum width=3cm, minimum height=0.9cm},
    arrow/.style={-Latex}
]

\node[box] (gm) {GameManager};

\node[box, below=of gm] (cm) {CircleManager};
\node[box, left=of cm] (chm) {CharacterManager};
\node[box, right=of cm] (dm) {DifficultyManager};

\node[box, below=of cm] (lg) {LevelGenerator};
\node[box, below=of lg] (rm) {RoomManager};
\node[box, below=of rm] (rg) {RoomGenerator};

\draw[arrow] (gm) -- (cm);
\draw[arrow] (gm) -- (chm);
\draw[arrow] (gm) -- (dm);

\draw[arrow] (cm) -- (lg);
\draw[arrow] (lg) -- (rm);
\draw[arrow] (rm) -- (rg);

\draw[arrow] (dm) -- (rg);

\end{tikzpicture}
\end{center}

\subsection{Gestión de personajes jugables}

El sistema contempla dos personajes jugables: Dante y Virgilio. Para garantizar extensibilidad, se ha diseñado un \textit{CharacterManager} responsable de:

\begin{itemize}
    \item Mantener el estado del personaje activo.
    \item Gestionar el posible cambio dinámico entre personajes.
    \item Centralizar atributos diferenciadores (estadísticas, habilidades, modificadores).
\end{itemize}

Ambos personajes comparten una interfaz común de comportamiento, lo que permite desacoplar la lógica del jugador del resto del sistema. De esta forma, la introducción de mecánicas específicas para cada personaje no afecta a los módulos de generación de niveles o enemigos.

\subsection{Modelo estructural del nivel}

La generación del nivel se basa en una arquitectura híbrida compuesta por:

\begin{itemize}
    \item Un modelo espacial basado en una cuadrícula bidimensional (grid) que controla la ocupación geométrica.
    \item Un modelo lógico basado en grafos donde cada nodo representa una habitación y cada arista una conexión navegable.
\end{itemize}

El grid se emplea exclusivamente para validar la coherencia espacial durante la generación. Cada habitación define un conjunto de celdas ocupadas (footprint), lo que permite soportar geometrías complejas como salas rectangulares o en forma de L.

Sin embargo, desde el punto de vista lógico, cada habitación constituye un único nodo en el grafo, independientemente de su tamaño. Esto permite desacoplar la estructura espacial de la conectividad del nivel.

\subsection{Proceso de generación procedural}

La generación se divide en dos fases claramente diferenciadas.

\subsubsection{Generación estructural}

En esta fase, el \textit{LevelGenerator} construye el grafo completo del círculo actual. El proceso sigue una estrategia de expansión iterativa:

\begin{enumerate}
    \item Se crea la sala inicial.
    \item Se selecciona una puerta libre de una sala existente.
    \item Se elige una plantilla de habitación compatible.
    \item Se verifica su encaje espacial en el grid.
    \item Si el encaje es válido, se actualizan grid y grafo.
\end{enumerate}

Una vez finalizada la generación, se determina la sala del jefe como el nodo más alejado de la sala inicial, utilizando un recorrido en anchura (BFS).

\subsubsection{Generación de contenido bajo demanda}

El contenido interno de las salas no se genera durante la fase estructural. En su lugar, el \textit{RoomGenerator} instancia enemigos, objetos y eventos únicamente cuando el jugador accede a una sala, así como en sus adyacentes inmediatas.

Este enfoque reduce el consumo de memoria y permite que el \textit{DifficultyManager} ajuste dinámicamente la complejidad del contenido en función del estado actual del jugador.

\subsection{Configuración basada en datos}

Cada uno de los nueve círculos se define mediante una estructura de configuración independiente que incluye:

\begin{itemize}
    \item Conjunto de enemigos disponibles.
    \item Conjunto de plantillas de habitaciones.
    \item Conjunto de objetos.
    \item Prefab del jefe.
    \item Parámetros de dificultad.
\end{itemize}

Esta aproximación sigue un paradigma \textit{data-driven}, separando configuración y comportamiento. Esto permite modificar el balance del juego o añadir nuevos círculos sin alterar la lógica central del sistema.

\subsection{Justificación arquitectónica}

Las decisiones adoptadas persiguen los siguientes objetivos:

\begin{itemize}
    \item Separación clara de responsabilidades.
    \item Bajo acoplamiento entre módulos.
    \item Escalabilidad para incorporar nuevas mecánicas.
    \item Soporte para geometrías complejas de habitaciones.
    \item Flexibilidad para integrar múltiples personajes jugables.
\end{itemize}

La combinación de un modelo espacial basado en grid y un modelo lógico basado en grafos permite mantener coherencia geométrica sin sacrificar expresividad estructural. Asimismo, la modularización del sistema facilita su mantenimiento y evolución futura.